\documentclass[11pt]{article}
\usepackage[letterpaper,margin=1in]{geometry}
\usepackage{amsmath,amssymb,amsthm,mathtools,booktabs,array,microtype}
\usepackage{xurl}
\usepackage[hidelinks]{hyperref}
\usepackage{xcolor}
\usepackage{enumitem}
\definecolor{ink}{RGB}{21,46,75}
\usepackage{titlesec}
\titleformat{\section}{\large\bfseries\color{ink}}{\thesection}{.65em}{}
\newtheorem{theorem}{Theorem}
\newtheorem{lemma}[theorem]{Lemma}
\newtheorem{corollary}[theorem]{Corollary}
\newtheorem{proposition}[theorem]{Proposition}
\theoremstyle{remark}\newtheorem{remark}[theorem]{Remark}
\newcommand{\Q}{\mathbb Q}
\newcommand{\C}{\mathbb C}
\newcommand{\R}{\mathbb R}
\newcommand{\ord}{\operatorname{ord}}
\newcommand{\Ind}{\operatorname{Ind}}
\newcommand{\Res}{\operatorname{Res}}
\newcommand{\Frob}{\operatorname{Frob}}
\newcommand{\Reg}{\operatorname{reg}}
\newcommand{\Fix}{\operatorname{Fix}}
\newcommand{\one}{\mathbf{1}}
\newcommand{\RH}{\mathrm{RH}}
\newcommand{\src}[1]{\href{https://github.com/lsernau/grh-lean/blob/72fead102f663d5c7a407964a13dc0cf9d977207/grh-compact/#1}{\nolinkurl{#1}}}
\hypersetup{pdftitle={Beyond Dirichlet GRH: Sylow-filtered prime races},pdfauthor={Research note}}
\title{\color{ink}\bfseries Beyond Dirichlet GRH:\\Sylow-filtered prime races and a two-pattern test\\for nonabelian zero exclusion}
\author{Research note and reproducible finite-group certificates}
\date{25 November 2026 (Pacific time)}
\begin{document}
\maketitle
\begin{abstract}
We isolate a route from the one-sided prime-error method in the compact Alp\"oge--Claude proof of Dirichlet GRH to the Riemann hypothesis for arbitrary Dedekind zeta functions. For a finite Galois extension $E/\Q$ with group $G$, let $H=[G,G]$ and choose one Sylow subgroup $P_\ell$ of $H$ for each prime $\ell\mid |H|$. We prove that, using the now established Dirichlet GRH, RH for $\zeta_E$ is equivalent to an upper bound of size $O_\varepsilon(x^{-1/2+\varepsilon})$ for a single centered, signed, smoothed rational-prime sum. Its weight is the sum of the characters $\Reg_G-\Ind_{P_\ell}^{G}\one$; only completely split primes and primes whose Frobenius has prime-power order in $H$ enter. The proof uses monomial representations of Sylow subgroups and a quantitative inequality on Heilbronn characters, not Artin holomorphy or auxiliary quadratic-twist selection. A second reduction for $G=\mathrm{PSL}_2(\mathbb F_7)$ uses only split primes and involution primes, with weights $21$ and $-1$. The corresponding entire $L$-function is a product of two quadratic Hecke $L$-functions over a field of degree 21. We prove the reductions, audit their interface with the Lean proof, and isolate the arithmetic estimate still missing. No new unconditional nonabelian RH is claimed.
\end{abstract}
\vspace{.4em}
\noindent\textbf{Status.} Theorems in this note are classical paper proofs. The upstream formal proof is independently machine-checked; this note has not been formalized. The bundled Python checks verify finite group computations, not analytic conjectures. Literature priority is not certified.
\setcounter{tocdepth}{1}
\tableofcontents
\newpage
\section{Scope, provenance, and what is different}
The public theorem statements in \src{src/Statements.lean} import only Mathlib. They assert Mathlib's RH and Dirichlet $L$-function nonvanishing to the right of $1/2$ for every positive modulus (except the principal pole), together with the corresponding zero-line statements. The compact project reports Lean 4.35.0, Mathlib commit \texttt{56a005855ea17e0b0958f73ec801ae140b0a1679}, and the three standard axioms \texttt{propext}, \texttt{Classical.choice}, \texttt{Quot.sound}. The external checks of Sutherland and Rubinstein report successful compilation, axiom checks, and replay of the compact project through the Lean kernel; further checks were in progress at their displayed update times \cite{sernau,sutherland,rubinstein}. We did not execute that build locally. The theorem is not a universal zero statement for Hecke, Artin, or automorphic $L$-functions.

We inspected the compact source at commit\newline\nolinkurl{72fead102f663d5c7a407964a13dc0cf9d977207}. The source map in Section~\ref{sec:bridge} identifies the exact declarations used to motivate the transfer. Nine prior ProveIt archives \cite{proveit} were retrieved through GitHub as base64 ZIP data, selected README files and available manuscript abstracts or section headings extracted. They already address, inter alia, class-group exponent six, CM coefficient fields, sparse singular-modulus relations, isogenies, arithmetic prime windows, prime bias, cubic Artin zero tests, and other one-sided criteria. See the manifest in the accompanying \texttt{prior-reports.md}. Their mathematical claims were not reverified here. The present note concentrates on a different organizing principle: use an \emph{entire product of induced Hecke functions whose zeros cannot conceal a nonabelian Dedekind zero}, with weights supported on a selected set of Frobenius classes. This also yields quantitative multiplicity bounds and a two-class example for the simple group of order 168.

\paragraph{What this establishes.} The reduction is an equivalence, not a proof of its arithmetic side. A completely split prime-counting estimate is already a familiar test for Dedekind RH; the benefit here is the ability to remove the abelian quotient, have zero logarithmic main term, discard composite-order and abelian-nonsplit Frobenius classes, and use a smaller Hecke factorization in particular examples. No claim is made that any of these advantages automatically makes the arithmetic estimate easier. The constituent tools (Artin formalism, Heilbronn characters, Landau's theorem, explicit formulas) are classical. The combined criteria and certificates are developed here; exhaustive priority checking has not been done.

\section{A one-sided analytic engine}
We work with a nontrivial entire product $F(s)$ of finitely many finite-order, nontrivial Hecke $L$-functions over fixed number fields. Positive integer powers are allowed. Regard it as an Euler product over rational primes, with all local factors included. Assume that the aggregate Euler coefficients are real, as happens for the representations below. Write
\[
D_F(s)=-\frac{F'(s)}{F(s)}=\sum_{p,k\ge1}\frac{a_F(p^k)\log p}{p^{ks}},\qquad \Re s>1,
\]
where $p$ runs over rational primes. By grouping the local ideal factors by rational norm, $|a_F(p^k)|\le C_F$ uniformly. For Artin representations the bound is simply the representation dimension. Such a product is nonzero in $\Re s\ge1$, has no pole, and $D_F(1)$ is finite and real. Define
\begin{equation}\label{eq:Rfull}
 R_F(x)=\sum_{p^k\le x}a_F(p^k)\log p\left(\frac1{p^k}-\frac1x\right)-D_F(1),\qquad x\ge1.
\end{equation}
The constant removes the main value at $1$. In contrast to a prime sum with a positive trivial-character coefficient, there is no logarithmic main term.

\begin{lemma}[One-sided zero test]\label{lem:analytic}
For $1/2\le\beta<1$, if $R_F(x)\le C x^{\beta-1}$ for all sufficiently large real $x$, then $F(s)$ is nonzero for $\Re s>\beta$. Conversely, if $F$ has no zeros in $\Re s>\beta$, then for every $\varepsilon>0$,
\[ R_F(x)=O_{F,\varepsilon}(x^{\beta-1+\varepsilon}).\]
In particular, zero exclusion to the right of $1/2$ is equivalent to the one-sided bound $R_F(x)\le C_{F,\varepsilon}x^{-1/2+\varepsilon}$ for every $\varepsilon>0$.
\end{lemma}
\begin{proof}
For $\Re s>1$, termwise integration (absolute convergence) yields
\begin{equation}\label{eq:mellin}
 T_F(s):=\int_1^\infty R_F(x)x^{-s}\,dx
 =\frac{D_F(s)/s-D_F(1)}{s-1}.
\end{equation}
Indeed the integral of $(1/n-1/x)\one_{x\ge n}$ is $n^{-s}/(s(s-1))$. The apparent singularity at 1 is removable. If $F$ has a zero of order $m>0$ at $\rho$ in $0<\Re\rho<1$, the residue of $T_F$ there is
\begin{equation}\label{eq:residue}
 -\frac{m}{\rho(\rho-1)}.
\end{equation}
In particular, at a \emph{real} zero in $(\beta,1)$ the residue is strictly positive. This sign is the central point: it is guaranteed by entireness, even if $F$ is built from a group with unknown Artin holomorphy.

After increasing $C$ and adding a bounded function $b(x)$ supported in a finite interval if necessary, $h(x)=C x^{\beta-1}-R_F(x)+b(x)\ge0$ for $x\ge1$. Its Mellin transform is initially convergent for real $s>1$ (a crude bound $R_F=O(1+\log^2 x)$ suffices). It continues meromorphically to $\Re s>\beta$ as
\[ U(s)=\frac C{s-\beta}-T_F(s)+\int_1^\infty b(x)x^{-s}\,dx.\]
Landau's lemma for a nonnegative transform says that if its abscissa of convergence $\sigma_c$ is finite, then the point $\sigma_c$ on the real axis is a singularity. For completeness, substitute $x=e^v$. If the transform extended holomorphically to a disk around $\sigma_c$, choose a real $s_0>\sigma_c$ close enough that its Taylor disk reaches $s_0-r<\sigma_c$. The Taylor series in the direction of decreasing $s$ has coefficients $\int_0^\infty h(e^v)e^{(1-s_0)v}v^j/j!\,dv\ge0$. Tonelli identifies its finite sum at $r$ with the integral at $s_0-r$, a contradiction. If $\sigma_c>\beta$, the only possible real singularity of $U$ there is a simple pole at a real zero of $F$ in $(\beta,1)$, with strictly negative residue by \eqref{eq:residue}. But $U(s)\ge0$ on the real axis to the right of $\sigma_c$, so such a pole is impossible. Thus $\sigma_c\le\beta$ (also allowing $-\infty$). The integral is holomorphic in $\Re s>\beta$; \eqref{eq:residue} excludes every zero of $F$ there, including nonreal ones.

For the converse, take $\sigma=\beta+\varepsilon<1$. The standard logarithmic-derivative estimate for fixed finite-order Hecke functions, from the Hadamard product of the completed function and the local zero-count bound, gives $D_F(s)=O_{F,\varepsilon}(\log(|\Im s|+3))$ in $\sigma\le\Re s\le2$ under the assumed zero exclusion. More explicitly, the sum of $1/(s-\rho)$ over zeros within unit height has $O_F(\log(|\Im s|+3))$ terms, each at distance at least $\varepsilon$; the remaining contribution including gamma factors has logarithmic size. The bound at small heights follows from compactness and nonvanishing. The smoothed Perron formula
\[
 S_F(x)=\frac1{2\pi i}\int_{(2)}D_F(s)\frac{x^{s-1}}{s(s-1)}\,ds
\]
represents the sum in \eqref{eq:Rfull} before centering: the inverse Mellin kernel is $(y-1)_+$, continuous at $y=1$. Shift to $\Re s=\sigma$; horizontal integrals tend to zero because of the $1/|s|^2$ decay. The only crossed pole is at $1$ with residue $D_F(1)$. The new vertical integral is $O_{F,\varepsilon}(x^{\sigma-1})$. For larger $\varepsilon$, the weaker assertion follows by shrinking it. Finally apply the first implication with $\beta=1/2+\varepsilon$ for all sufficiently small $\varepsilon$ to get the equivalence.
\end{proof}

\paragraph{Prime-only form.} Let $S$ be any fixed finite set of rational primes containing the ramified primes, and suppose $a_F(p)=w(\Frob_p)$ outside $S$. Put
\begin{align}
 C_{F,S}&=D_F(1)-\sum_{\substack{p,k\ge1:\ p\in S\ \text{or}\ k\ge2}}\frac{a_F(p^k)\log p}{p^k},\label{eq:C}\\
 R_w(x)&=\sum_{\substack{p\le x\\p\notin S}}w(\Frob_p)\log p\left(\frac1p-\frac1x\right)-C_{F,S}.
 \label{eq:Rprime}
\end{align}
The correction series in \eqref{eq:C} is absolutely convergent. The omitted coefficients with $p^k\le y$ have total absolute log-weight $O_F(\sqrt y\log^2(2y))$ (sum separately over $k=2$ and $3\le k\le \log_2 y$, plus finitely many prime powers at $S$). Partial summation then bounds the omitted tail and the $1/x$ contribution by $O_F(x^{-1/2}\log^2(2x))$. Consequently
\begin{equation}\label{eq:correction}
 R_w(x)=R_F(x)+O_{F,S}(x^{-1/2}\log^2(2x)).
\end{equation}
The equivalence in Lemma~\ref{lem:analytic} is unchanged at the $1/2+\varepsilon$ scale. The constant can also be defined independently of any RH assumption as the limit of the smoothed prime sum: the fixed-conductor Hecke prime number theorem gives $S_F(x)\to D_F(1)$, and the omitted sum tends to the correction in \eqref{eq:C}. This is a definition of the exact constant, not a finite numerical estimate. For rational rescalings of $w$ by a positive factor, rescale the constant as well.

\section{A general nonabelian reduction}\label{sec:sylow}
Let $E/\Q$ be finite Galois, with group $G$. Let $H\le G$ be any subgroup and put $B=E^H$. We will assume RH for $\zeta_B$, known from Dirichlet GRH when $B/\Q$ is abelian. In particular, take $H=G'=[G,G]$, so that $B$ is the maximal abelian subextension of $E$. If $H=1$, $E=B$ and the target is already covered; assume $H\ne1$. For each prime $\ell$ dividing $|H|$, choose a Sylow $\ell$-subgroup $P_\ell\le H$. Define the honest representation character
\begin{equation}\label{eq:W}
 W_H=\sum_{\ell\mid|H|}\Ind_{P_\ell}^G(\Reg_{P_\ell}-\one)
 =\sum_{\ell\mid|H|}(\Reg_G-\Ind_{P_\ell}^G\one).
\end{equation}
Every summand has trivial-character multiplicity zero. Artin formalism gives
\begin{equation}\label{eq:prod}
 F_H(s)=L(s,W_H,E/\Q)=\prod_{\ell\mid|H|}\frac{\zeta_E(s)}{\zeta_{E^{P_\ell}}(s)}.
\end{equation}
Each factor is entire by Aramata--Brauer. More usefully for this note, because $P_\ell$ is a finite $\ell$-group, it is monomial (every irreducible is induced from a linear character of a subgroup). For every nontrivial irreducible the inducing character must be nontrivial, since an induction of the trivial character contains the trivial representation. Thus every factor in \eqref{eq:prod}, on decomposing the augmentation representation, is a product of nontrivial finite-order Hecke $L$-functions with positive powers. This directly supplies the analytic hypotheses of Lemma~\ref{lem:analytic}. No claim about holomorphy for the individual nonmonomial Artin irreducibles of $G$ has been used.

\begin{theorem}[Sylow-filtered prime race]\label{thm:sylow}
Assume RH for $\zeta_B$. For $F_H$ and $W_H$ as in \eqref{eq:W}--\eqref{eq:prod}, with $R_{W_H}$ centered by \eqref{eq:C}, the following are equivalent:
\begin{enumerate}[label=(\roman*)]
\item RH for $\zeta_E$;
\item $F_H$ has no zeros in $\Re s>1/2$;
\item For every $\varepsilon>0$, $R_{W_H}(x)\le C_{E,H,\varepsilon}x^{-1/2+\varepsilon}$ for all sufficiently large $x$.
\end{enumerate}
If these hold, every irreducible Artin $L$-function of $E/\Q$ is holomorphic and nonzero in $\Re s>1/2$, except the trivial pole at $1$. This last assertion is local zero/pole exclusion in the open half-plane, not global Artin holomorphy.
\end{theorem}

We first prove a stronger local estimate that avoids any unjustified nonnegativity of Artin zero orders. Fix $\rho\ne1$ in $0<\Re\rho<1$. Let $n_\chi=\ord_{s=\rho} L(s,\chi,E/\Q)$ for $\chi\in\operatorname{Irr}(G)$, where poles have negative order. Brauer induction gives meromorphic continuation, so these are integers. The Heilbronn virtual character is
\[ \Theta=\sum_\chi n_\chi\chi,\qquad d=\Theta(1)=\ord_\rho\zeta_E.\]
By Artin formalism and Frobenius reciprocity, for a subgroup $K\le G$ and an irreducible character $\eta$ of $K$,
\begin{equation}\label{eq:heil}
 \langle\Res_K\Theta,\eta\rangle=\ord_\rho L(s,\eta,E/E^K).
\end{equation}
For cyclic $K$ or for a prime-power group $K$, the right side is nonnegative: every irreducible is monomial and its Hecke function is holomorphic at $\rho$, including the trivial case. Thus $\Res_K\Theta$ is an actual character (possibly zero). In particular $d\ge0$; $d=0$ implies $\Theta=0$, since a genuine character of dimension zero on every cyclic subgroup is identically zero. These are classical Heilbronn--Stark facts \cite{heilbronn}.

\begin{proposition}[Multiplicity domination]\label{prop:dom}
For each $\ell\mid|H|$, let $e_\ell=\ord_\rho(\zeta_E/\zeta_{E^{P_\ell}})$, and $b=\ord_\rho\zeta_B$. Then $0\le e_\ell\le d$ and
\begin{equation}\label{eq:dom}
 d-b\le 2\sum_{\ell\mid|H|}\delta_\ell e_\ell\le 2\ord_\rho F_H,
 \qquad \delta_\ell=\frac{\#\{h\in H:\ell\mid\ord(h)\}}{|H|}.
\end{equation}
\end{proposition}
\begin{proof}
The genuine character $\Res_{P_\ell}\Theta$ has dimension $d$ and its trivial multiplicity is $\ord_\rho\zeta_{E^{P_\ell}}$ by \eqref{eq:heil}. Therefore $e_\ell$ is the codimension of its fixed subspace. For $h\in H$, write $h=\prod_{\ell\mid\ord(h)} h_\ell$ for the commuting prime-power components of $h$ in the cyclic group $\langle h\rangle$. Each $h_\ell$ is conjugate within $H$ into $P_\ell$. In the genuine representation associated to $\Res_{\langle h\rangle}\Theta$, the fixed subspace of $h_\ell$ has codimension at most $e_\ell$: its dimension agrees by restriction of characters with the fixed-space dimension for the element in $P_\ell$. The fixed subspace of $h$ is the intersection of those of the $h_\ell$, because each component is a power of $h$. Thus its codimension is at most $\sum_{\ell\mid\ord(h)}e_\ell$. Eigenvalues of $h$ have modulus one. Each nonfixed eigenvalue contributes at most 2 to $d-\Re\Theta(h)$, giving
\[
 0\le d-\Re\Theta(h)\le2\sum_{\ell\mid\ord(h)}e_\ell.
\]
Average over $h\in H$. By \eqref{eq:heil}, the average of $\Theta(h)$ equals the trivial multiplicity $b$, whether or not $\Res_H\Theta$ itself is genuine. This proves \eqref{eq:dom}; the sum of the $e_\ell$ is the order of the product \eqref{eq:prod}.
\end{proof}
\begin{proof}[Proof of Theorem~\ref{thm:sylow}]
At any $\rho$ in $1/2<\Re\rho<1$, RH for $B$ says $b=0$. Proposition~\ref{prop:dom} gives
\begin{equation}\label{eq:mult}
 \tfrac12\ord_\rho\zeta_E\le\ord_\rho F_H\le r\ord_\rho\zeta_E,\qquad r=\#\{\ell:\ell\mid|H|\}.
\end{equation}
Thus $F_H$ and $\zeta_E$ have exactly the same zero support in that half-strip. They are nonzero on $\Re s\ge1$ apart from the pole of $\zeta_E$; the functional equation of $\zeta_E$ supplies the other half of the critical strip. The equivalence with (iii) follows from Lemma~\ref{lem:analytic} and \eqref{eq:correction}. Finally $d=0$ forces $\Theta=0$ by cyclic restriction. At $s=1$, every nontrivial Artin factor is nonzero and holomorphic by the standard order calculation using Brauer induction and the simple pole of each trivial Hecke function. Indeed each induced linear character contributes order $-1$ precisely when the inducing character is trivial, which equals the negative of its trivial multiplicity; by linearity the order for any irreducible other than the trivial one is zero. For $\Re s\ge1$, $s\ne1$, induction from Hecke factors gives nonzero and holomorphic. This proves the last assertion.
\end{proof}

\subsection{What primes enter?}
For a subgroup $P\le G$, $\Ind_P^G\one(g)=\#\Fix(g,G/P)=|\{x\in G:x^{-1}gx\in P\}|/|P|$. Hence
\begin{equation}\label{eq:weights}
 W_H(g)=\begin{cases}
 r|G|-\displaystyle\sum_\ell [G:P_\ell],&g=1,\\
 -\displaystyle\sum_\ell \#\Fix(g,G/P_\ell),&g\ne1.
 \end{cases}
\end{equation}
For $g\ne1$ a term can be nonzero only if $g$ is conjugate to a prime-power-order element of $H$. When $H$ is normal, as for $G'$, this says $g\in H$ and $g$ has prime-power order. The prime race therefore ignores every prime not splitting completely in $B$, and every nonidentity composite-order Frobenius class. For $H=G'$, this is exactly where the established Dirichlet GRH is used: the abelian quotient is already controlled. The weight has average zero over $G$, either from the trivial multiplicity or by the orbit-count formula for a transitive coset action. One may add a constant to make all weights nonnegative; its centered prime error is already of the required size by RH for $\zeta_{\Q}$, so this operation leaves the criterion unchanged. This can be useful if an integrality construction requires nonnegative exponents, but it does not produce that construction.

\begin{corollary}[Arbitrary number fields]\label{cor:any}
Let $K$ be any number field and $E$ its Galois closure over $\Q$. Establishing the single estimate (iii) for $E$, with $H=G'$, proves RH for $\zeta_K$.
\end{corollary}
\begin{proof}
$E/K$ is Galois, so $\zeta_E/\zeta_K$ is entire by Aramata--Brauer, or use \eqref{eq:heil} with the trivial character of $\operatorname{Gal}(E/K)$. A zero of $\zeta_K$ in the right half-strip would be a zero of $\zeta_E$. The functional equation completes the conclusion.
\end{proof}

\section{Explicit symmetric-group tests}
The following weights are primitive integer rescalings of \eqref{eq:weights}, with $H=A_n$ for $G=S_n$. They concern unramified rational primes, classified by the cycle type of Frobenius on the roots of a polynomial with the indicated Galois group. They can be evaluated by factoring the polynomial modulo $p$, omitting primes dividing its discriminant. Weights not shown are zero.
\begin{center}
\begin{tabular}{llll}
\toprule
Group & Sylow orders in $H$ & Cycle types & Weights\\
\midrule
$S_4$ & $4,3$ & $1^4;\;2^2;\;1\,3$ & $17,-3,-1$\\
$S_5$ & $4,3,5$ & $1^5;\;1\,2^2;\;1^2\,3;\;5$ & $133,-3,-2,-2$\\
\bottomrule
\end{tabular}
\end{center}
For example, in $S_5$ the Sylow subgroups of $A_5$ may be a Klein four group, a cyclic group of order 3, and a cyclic group of order 5. Their indices in $S_5$ are $30,40,24$. At the identity, \eqref{eq:weights} is $360-94=266$. A double transposition has centralizer of size 8 and the Klein four contains three of them, giving weight $-8\cdot3/4=-6$. A 3-cycle has centralizer of size 6 and the order-three subgroup contains two, giving $-4$. A 5-cycle has centralizer of size 5 and the order-five subgroup contains four, giving $-4$. Divide by 2. The $S_4$ calculation is similar: the indices are $6,8$, and the nonidentity values are $-6,-2$. The full enumeration in \texttt{verify\_groups.py} verifies these results directly from the coset formula, including zero weights.

\paragraph{Example polynomial.} $f(X)=X^5-X-1$ has discriminant $2869=19\cdot151$. It is irreducible modulo 3; modulo 2 it has type $(2,3)$. The finite factorizations are checked in the accompanying script. Irreducibility over $\Q$ implies a transitive group of prime degree 5, hence primitive. The cube of a $(2,3)$ element is a transposition. A primitive permutation group containing a transposition is the full symmetric group: the connected components of the graph of conjugate transpositions are a block system, so the graph is connected; edge transpositions of a connected graph generate the full symmetric group. Therefore this is an explicit $S_5$ instance, not a hypothetical family. The established Dirichlet GRH handles its quadratic abelian subfield $\Q(\sqrt{2869})$.

The resulting arithmetic challenge is a signed prime race, not an ordinary congruence-class count:
\begin{equation}\label{eq:s5}
 \begin{split}
 &\sum_{p\le x,\ p\nmid2869}\! \log p\,(1/p-1/x)\,
 \big(133\one_{1^5}-3\one_{1\,2^2}-2\one_{1^2\,3}-2\one_{5}\big)\\
 &\hspace{3cm}\le C_{S_5}+O_\varepsilon(x^{-1/2+\varepsilon}).
 \end{split}
\end{equation}
Here the indicators refer to factor degrees of $f\bmod p$ and $C_{S_5}$ is the rescaled constant \eqref{eq:C}. For every $\varepsilon>0$, the bound is equivalent to Dedekind RH for the splitting field (not just the degree-five field). This is a sufficient but still unproved new arithmetic target.

\section{A sharper two-pattern reduction for \texorpdfstring{$\mathrm{PSL}_2(7)$}{PSL2(7)}}\label{sec:psl}
Consider $G=\mathrm{GL}_3(\mathbb F_2)\simeq\mathrm{PSL}_2(\mathbb F_7)$, simple of order 168. Its conjugacy class orders are $1,2,3,4,7,7$. We denote the six irreducible characters by $1,\chi_3,\overline\chi_3,\chi_6,\chi_7,\chi_8$. The standard character table \cite{groupnames} is
\begin{center}
\begin{tabular}{crrrrrr}
\toprule
 & $1$ & $2$ & $3$ & $4$ & $7A$ & $7B$\\
Class size & 1 & 21 & 56 & 42 & 24 & 24\\
\midrule
$1$ & 1 & 1 & 1 & 1 & 1 & 1\\
$\chi_3$ & 3 & -1 & 0 & 1 & $a$ & $\bar a$\\
$\overline\chi_3$ & 3 & -1 & 0 & 1 & $\bar a$ & $a$\\
$\chi_6$ & 6 & 2 & 0 & 0 & -1 & -1\\
$\chi_7$ & 7 & -1 & 1 & -1 & 0 & 0\\
$\chi_8$ & 8 & 0 & -1 & 0 & 1 & 1\\
\bottomrule
\end{tabular}
\qquad $a=(-1+\sqrt{-7})/2$.
\end{center}
Let $P=D_8=\langle r,s:r^4=s^2=1,\ srs=r^{-1}\rangle$ be a Sylow 2-subgroup. Define two quadratic characters of $P$ by
\[ \lambda_A(r)=1,\quad\lambda_A(s)=-1;\qquad
 \lambda_B(r)=-1,\quad\lambda_B(s)=1.\]
Let $J\simeq S_4$ be a stabilizer of a nonzero vector in $\mathbb F_2^3$, and $U\simeq C_7\rtimes C_3$ the normalizer of a Sylow 7-subgroup. Let $\omega$ be a nontrivial order-three character of $U/C_7$. The following induction identities follow by taking inner products with the displayed table:
\begin{align}
 \Ind_P^G\lambda_A&=\chi_3+\overline\chi_3+\chi_7+\chi_8,&
 \Ind_P^G\lambda_B&=\chi_6+\chi_7+\chi_8,\label{eq:indA}\\
 \Ind_J^G\one&=\one+\chi_6,&
 \Ind_J^G\operatorname{sgn}&=\chi_7,&
 \Ind_U^G\omega&=\chi_8.\label{eq:indB}
\end{align}
To verify these identities without a character algebra package, write $c_j=\chi(g_j)$ for a character of $G$ on a class of order $j$. The inner products for the first two inductions are $(c_1-3c_2+2c_4)/8$ and $(c_1+c_2-2c_4)/8$. Indeed $P$ contains five involutions and two elements of order four; the weighted sums of the two linear characters on these classes are $(-3,2)$ and $(1,-2)$. For $J$, the ordinary inner product is $(c_1+9c_2+8c_3+6c_4)/24$ and its sign version is $(c_1-3c_2+8c_3-6c_4)/24$, by the cycle counts in $S_4$. For $U$, the 14 elements outside $C_7$ contribute $7(\omega+\bar\omega)c_3=-7c_3$, and its six elements of order seven contribute $3c_{7A}+3c_{7B}$. Divide the resulting sum by 21. The companion Python program enumerates the group as binary invertible matrices, verifies the subgroup orders and signed counts, and checks all the multiplicities by exact rational arithmetic.

Set $F_P=L(s,\Ind_P^G\lambda_A)L(s,\Ind_P^G\lambda_B)$. This is a product of two quadratic Hecke functions over $E^P$ (degree 21 over $\Q$). Its combined character is
\begin{equation}\label{eq:pslweight}
 W_P=(42,-2,0,0,0,0)
\end{equation}
on the six classes in table order. Thus its prime weight, divided by 2, is $21$ at a split prime, $-1$ at an involution prime, and zero elsewhere. These two tests are substantially more selective than the general Sylow race for this group.

\begin{theorem}[Two-pattern septic test]\label{thm:psl}
Assume RH for $\zeta_{\Q}$. For every Galois extension $E/\Q$ with $G\simeq\mathrm{PSL}_2(7)$, RH for $\zeta_E$ is equivalent to
\begin{equation}\label{eq:psltest}
 21\sum_{\substack{p\le x\\\Frob_p=1}}\log p(1/p-1/x)
 -\sum_{\substack{p\le x\\\ord(\Frob_p)=2}}\log p(1/p-1/x)
 -C_P\ \le\ O_{E,\varepsilon}(x^{-1/2+\varepsilon})
\end{equation}
for every $\varepsilon>0$, omitting a fixed finite set containing all ramified primes. The constant $C_P$ is half the constant \eqref{eq:C} for $F_P$ and that omitted set. The test does not presuppose Artin holomorphy. Furthermore, at any $\rho$ in $1/2<\Re\rho<1$, $\ord_\rho\zeta_E\le 6\ord_\rho F_P$.
\end{theorem}
\begin{proof}
At such a $\rho$, write $n_1,n_3,n_{\bar3},n_6,n_7,n_8$ for Artin orders, $u=n_3+n_{\bar3}$, and $d=\ord_\rho\zeta_E$. The known RH gives $n_1=0$. By \eqref{eq:indB}, $n_6=\ord_\rho\zeta_{E^J}\ge0$, $n_7=\ord_\rho L(s,\operatorname{sgn},E/E^J)\ge0$, and $n_8=\ord_\rho L(s,\omega,E/E^U)\ge0$. We do not assume that $n_3$ or $n_{\bar3}$ is nonnegative.

Let $a_P,b_P$ be the two Hecke zero orders in \eqref{eq:indA}; both are nonnegative and satisfy
\[ a_P=u+n_7+n_8,\qquad b_P=n_6+n_7+n_8.\]
By taking degrees in the regular character,
\begin{align*}
 d&=3u+6n_6+7n_7+8n_8\\
  &=3a_P+6b_P-2n_7-n_8\le 6(a_P+b_P).
\end{align*}
If $F_P$ is zero-free at $\rho$, then $a_P+b_P=0$ and $d=0$. Conversely $d=0$ implies $\Theta=0$ by cyclic restriction as above, so $a_P=b_P=0$. Thus $F_P$ and $\zeta_E$ have the same right-half-strip zero support. Apply Lemma~\ref{lem:analytic}, \eqref{eq:correction}, the factor $1/2$ in \eqref{eq:pslweight}, and the functional equation of $\zeta_E$.
\end{proof}

\paragraph{A concrete polynomial.} Trinks' polynomial $g(X)=X^7-7X+3$ has Galois group $\mathrm{PSL}_2(7)$, as documented by Erbach, Fischer and McKay \cite{trinks}. We use this published group determination, not merely a few sample factorizations. Its discriminant is $3^8 7^8$. In the seven-point action the cycle types are $1^7$, $1^3 2^2$, $1\,3^2$, $1\,2\,4$, and $7$. The group enumeration checks this list. In \eqref{eq:psltest}, the first sum is therefore over primes for which $g$ splits into seven linear factors and the second over primes with three linear and two irreducible quadratic factors. Omit 3 and 7. All other primes, including irreducible septic reductions, can be ignored. In fact the weight $W_P$ is half the induction of the nontrivial character of an order-two subgroup; the factorization \eqref{eq:indA} reduces the Hecke base degree from 84 to 21 without changing the zero test. The certificate concerns identities in the finite group, not the infinite prime inequality.

\section{Transfer from the Lean proof: where the hard work is}\label{sec:bridge}
\subsection{Source-level map}
The main identity used in the compact proof is a comparison of upper and lower estimates on a common squared covolume of a family of theta lifts. We checked the following source declarations, not merely their names. Write $m$ for the number of lifts, $a$ for the long order parameter, $f$ for the number of infinite places, $b$ for the large scale, $h$ for the number of thick places and $h_0$ for the number of top places. Put $\mathcal B=fa(\log b)^3$.
\begin{center}
\small
\begin{tabular}{@{}p{.44\linewidth}p{.51\linewidth}@{}}
\toprule
Source/declaration & Role and limitation\\
\midrule
\src{src/New/Endgame/ShortAndEnough.lean}\newline \texttt{Data.short\_forms}, \texttt{Data.enough\_forms}&Construct many short integral forms with small index and almost full dimension. Specific unitary groups and theta-lift data.\\[2pt]
\src{src/New/Endgame/LongForms.lean}\newline \texttt{Data.integral\_forms\_long}&Lower covolume bound via fit twists, lift-volume ratios and height estimates. Not a theorem for arbitrary prime weights.\\[2pt]
\src{src/New/Main.lean}\newline \texttt{star\_inequality}&Cancels the common height term, giving $m\,\mathrm{rem}\le(\log b)^2\mathcal B$.\\[2pt]
\src{src/New/Endgame/ShortAndEnough.lean}\newline \texttt{Data.remainder\_eq}&Exact finite prime-sum identity: $\mathrm{rem}=a(h(R_M(x)-R_M(x_1))+h_0 R_1(x_1))$. Relies on cyclotomic place-class count.\\[2pt]
\src{src/New/Endgame/PrimeError.lean}\newline \texttt{primeError}, \texttt{isPolar\_primeError}&Weighted non-split Dirichlet prime-power sum, its transform and residues.\\[2pt]
\src{src/New/Endgame/OneSidedCriterion.lean}\newline \texttt{IsPolar.exists\_window\_sub\_mul\_le}&Generic auxiliary-scale selection using moment estimates.\\[2pt]
\src{src/Lib/Analytic/Landau/PolarLandau.lean}\newline \texttt{IsPolar.landau}&One-sided exclusion from sign of real residues. Transferable analytic mechanism.\\[2pt]
\src{src/Lib/Analytic/QuadraticTwist/ZeroFreeTwist.lean}\newline \texttt{exists\_entireL\_quadTwist\_ne\_zero}&Dirichlet quadratic twist separation. Avoided in the proposed Hecke-product tests.\\
\bottomrule
\end{tabular}
\end{center}
The actual prime weight in \texttt{PrimeError.lean} is $\Lambda(n)\one_{n\equiv1\,\bmod M}(1-\chi_K(n))$. Its transformed residues are differences of zero multiplicities of characters and their quadratic twists. The proof needs the twist-selection lemma so a target zero contributes with the required sign. Our $F_H$ and $F_P$ already have nonnegative orders; no twist family or unproved zero-density theorem for nonabelian twists is required in the analytic stage. The aggregate weights are real and signed. Adding a constant can render them nonnegative, at the cost of a known Dirichlet main term and error, but does not identify the source's jet valuations with those weights.

\subsection{A sufficient replacement for the main inequality}
For either of our target prime errors, let $R=R_F$ be the full prime-power error \eqref{eq:Rfull} and let $R_0$ be a known RH/GRH error (for example the smoothed rational-prime-power error of $\zeta_{\Q}$ with its logarithmic main term and constant removed). In analogy with the source, it would suffice, for each small $\varepsilon>0$ and every sufficiently large $x$, to construct data with adjustable $x_1$ satisfying
\begin{align}
 & m\,a h\,[R(x)-R(x_1)+X R_0(x_1)]\le fa(\log b)^5,\label{eq:star-new}\\
 &m\ge\sqrt x,\qquad f/h\ll\log b,\qquad \log b=x^\varepsilon,\qquad
 0\le X\le x^A,\nonumber
\end{align}
for a fixed $A$ and for all $\log x_1\in[t,t^2]$ with $t=C\log x$ and fixed $C$ sufficiently large (the constants may depend on the fixed field). The analogue of the source's window lemma applies at the classical analytic level to $v\mapsto R(e^v),R_0(e^v)$: they have exponential growth bounds and their transforms are holomorphic at and to the right of 1, with the poles of the logarithmic derivative simple. Alternatively, for the fixed Hecke product, the standard unconditional fixed-conductor zero-free region and prime ideal theorem give $R(x_1)\ll\exp(-c\sqrt{\log x_1})$ after possibly changing $c>0$. Fixed exceptional zeros can be absorbed in this asymptotic bound. Choose $\log x_1=t^2$ with $C$ large; then $R(x_1)-XR_0(x_1)=O(x^{-2})$ if the base error has its known square-root bound. Dividing \eqref{eq:star-new} gives
\[
 R(x)\le O\left(\frac{f\log^5 b}{mh}\right)+O(x^{-2})
      =O(x^{-1/2+6\varepsilon}).
\]
The arbitrary $\varepsilon$ and Lemma~\ref{lem:analytic} would conclude RH for the target. This is a complete implication, but \eqref{eq:star-new} has not been proved for these weights. Requiring it on the full window makes the scale selection explicit, rather than silently moving a datum whose arithmetic construction may depend on $x_1$.

\subsection{A rigorous obstruction to a naive substitution}
\begin{proposition}\label{prop:obstruction}
For $G'\ne1$ and $H=G'$, the weight $W_H(\Frob_p)$ cannot agree on all sufficiently large unramified primes with any function of a fixed residue class $p\bmod M$. The same holds for $W_P$ of \eqref{eq:pslweight}.
\end{proposition}
\begin{proof}
Let $J=E\cap\Q(\zeta_M)$. Since $J/\Q$ is abelian, every $g\in G'$ acts trivially on $J$. In the Galois group of the compositum $E\Q(\zeta_M)$, the pairs $(1,1)$ and $(g,1)$ therefore exist. Chebotarev applied to the two conjugacy classes gives infinitely many rational primes $p\equiv1\pmod M$ with these respective Frobenius data. Pick any nonidentity element $g$ in a Sylow subgroup of $G'$. The two values of $W_H$ differ: the value at 1 is positive, the other is strictly negative by the coset formula. In the simple group case take an involution. This contradicts dependence on a fixed residue class alone.
\end{proof}
This is why the transfer is not an application of existing Dirichlet GRH estimates. The argument of \texttt{Data.remainder\_eq} uses a cyclotomic place count in an essential way. A nonabelian counterpart would require a new carrier of Frobenius information and a new proof that its integrality exponents realize the weighted coset fixed-point function. The existing upper/lower height estimates also have to survive that change. It is not enough to replace a Dirichlet symbol by a character table entry in the final line of the Lean proof.

\subsection{Concrete milestones for a successor proof}
\begin{enumerate}[leftmargin=*,label=\arabic*.]
\item Start with the septic example. Seek a nonnegative integral local exponent whose prime contribution realizes a positive translate of $W_P/2$. The test has only two nonzero Frobenius classes before translation. This is an exact target at all prime powers after adding the harmless correction \eqref{eq:correction}.
\item Work over the degree-21 field and its two explicit quadratic extensions, rather than over the entire degree-168 field where possible. Determine whether a theta-lift/jet construction can encode both quadratic Hecke characters with uniform local control. The factorization \eqref{eq:indA} is a simplification of the \emph{target} $L$-functions, not of the existing automorphic construction.
\item Establish the counterparts of A (short forms), B (enough forms), C (integral forms long), and the exact remainder identity. Audit the cancellation of the same main term and prove a bound uniform in the auxiliary scale. The single critical numerical inequality is \eqref{eq:star-new}.
\item Apply the analytic theorem and the character certificates. If the arithmetic step is proved, no additional quadratic-twist nonvanishing or Artin holomorphy input is needed. If a different construction yields the bound only for a fixed $\varepsilon>0$, the argument gives zero exclusion only to the corresponding line, not full RH.
\end{enumerate}

\section{What other conjectures might benefit?}
\begin{center}
\small
\begin{tabular}{@{}p{.25\linewidth}p{.40\linewidth}p{.28\linewidth}@{}}
\toprule
Target & Specific mechanism & Current status in this report\\
\midrule
Dedekind RH, nonabelian Galois closure & Sylow-filtered race, or sharper monomial detection in small groups & Proved equivalent criterion; bound open\\
Artin nonvanishing / holomorphy off critical line & Vanishing of $\Theta(1)$ implies zero Heilbronn character & Follows if preceding bound proved; no global Artin holomorphy result\\
Finite-order Hecke GRH over nonabelian fields & Include their abelian extension in a Galois closure; its zeta factorization and holomorphy transfer the zero exclusion & Same missing arithmetic step\\
Automorphic $L$-functions with infinite image & One-sided analytic criterion for an entire real Euler product may generalize given appropriate coefficient bounds; no finite Frobenius carrier & No arithmetic transfer established\\
BSD, Bloch--Kato, general Artin holomorphy on the line & Congruence/height modules suggest methods, but our nonvanishing result controls neither central zero orders nor algebraic ranks & No proof or reduction claimed\\
Twin primes / binary Goldbach & Correlations between different integer arguments, not one-variable Frobenius distributions & No transfer established\\
\bottomrule
\end{tabular}
\end{center}
The finite-order Hecke assertion can be spelled out: a character of finite order over a number field $K$ corresponds to an abelian extension $L/K$ by class field theory. If the Galois closure $E/\Q$ of $L$ satisfies RH, then RH for $\zeta_L$ follows from Corollary~\ref{cor:any}; the factorization $\zeta_L=\prod_{\chi}L(s,\chi)$ over $K$ consists of holomorphic factors in the right half-strip, hence every finite-order Hecke factor is zero-free there. The procedure could therefore in principle reach the full finite-order Hecke GRH, not just the Dirichlet case. It says nothing about general infinite-order Hecke functions or arbitrary automorphic $L$-functions without further input.

\section{Reproducibility and audit}
The included \texttt{verify\_groups.py} enumerates $S_4$, $S_5$, and $\mathrm{GL}_3(2)$, computes their weights by counting conjugation into subgroups, and verifies the induction multiplicities in \eqref{eq:indA}--\eqref{eq:indB} using exact fractions. The published character table is an input, not inferred from floating point. Its output is \texttt{group\_certificate.json}. The optional \texttt{sample\_prime\_races.py} uses SymPy to check polynomial discriminants and sample splitting types, and evaluates uncentered finite prime sums. Those samples are illustrative only and do not prove asymptotic estimates. The upstream identity table in \texttt{source\_map.tsv} gives Git blob hashes for the inspected files; \texttt{audit\_source.py} checks these against a separately supplied Git checkout. There is no network-dependent computation in the group certificate. The archive contains source, compiled manuscript, scripts, data and SHA-256 checksums.

\paragraph{Novelty and limitations.} The universal proof uses a classical virtual-character method in a new-to-this-note combination with the upstream smoothing. The main additions relative to the prior reports whose summaries we read are a single Sylow-filtered race valid for arbitrary finite Galois groups, the quantitative order inequality \eqref{eq:dom}, and an explicit two-pattern certificate for the order-168 group. This is not a claim that equivalent formulations are absent from the literature. We did not review all prior ZIP members, run the upstream Lean build, or prove a nonabelian covolume inequality. No previously open conjecture is certified solved by the new reductions themselves. The report's contribution is to replace a broad target by precise sufficient arithmetic inequalities and remove an auxiliary twist and Artin-holomorphy obstacle.

\begin{thebibliography}{99}\small
\bibitem{sernau} L. Sernau and collaborators, \emph{GRH for Dirichlet L-functions, in Lean 4}, 25 Nov. 2026. \href{https://github.com/lsernau/grh-lean/tree/72fead102f663d5c7a407964a13dc0cf9d977207}{Pinned repository}. Original Alp\"oge--Claude proof: \url{https://alpo.ge/remark.pdf}.
\bibitem{sutherland} A. Sutherland, \emph{Verification files for the shorter Lean proof of RH and GRH}, \url{https://math.mit.edu/~drew/shortgrh/}, status page consulted 25 Nov. 2026 Pacific time (displayed update 26 Nov. UTC).
\bibitem{rubinstein} M. Rubinstein, \emph{An independent check of the short Lean proof of RH and GRH}, \url{https://mrubinst.github.io/grh-lean-compact-check/}, status page consulted 25 Nov. 2026 Pacific time.
\bibitem{heilbronn} R. Foote, H. Ginsberg, V. Kumar Murty, \emph{Heilbronn characters}, Bull. Amer. Math. Soc. 52 (2015), 465--496. \href{https://doi.org/10.1090/BULL/1492}{doi:10.1090/BULL/1492}. Classical restriction, Artin orders and monomial subgroup facts.
\bibitem{neukirch} J. Neukirch, \emph{Algebraic Number Theory}, Springer, 1999, Ch. VII (Hecke and Artin $L$-functions, Artin formalism).
\bibitem{mv} H. L. Montgomery, R. C. Vaughan, \emph{Multiplicative Number Theory I: Classical Theory}, Cambridge University Press, 2007, Ch. 15 (Landau's lemma and sign changes).
\bibitem{ik} H. Iwaniec, E. Kowalski, \emph{Analytic Number Theory}, AMS Colloquium Publications 53, 2004, Ch. 5 (completed $L$-functions and explicit formulas). Fixed-field prime ideal estimates are standard classical inputs.
\bibitem{groupnames} T. Dokchitser, GroupNames, \emph{Character tables}, entry $\mathrm{GL}_3(\mathbb F_2)$, \url{https://www.people.maths.bris.ac.uk/~matyd/GroupNames/characters.html}.
\bibitem{trinks} D. W. Erbach, J. Fischer, J. McKay, \emph{Polynomials with PSL(2,7) as Galois group}, J. Number Theory 11 (1979), 69--75. \href{https://doi.org/10.1016/0022-314X(79)90020-9}{doi:10.1016/0022-314X(79)90020-9}. Records Trinks' polynomial $X^7-7X+3$.
\bibitem{proveit} ProveIt, \href{https://github.com/VladimirReshetnikov/ProveIt/tree/main/docs/incoming}{Incoming report directory}, retrieved 25 Nov. 2026; nine archives sampled as listed in \texttt{prior-reports.md}.
\end{thebibliography}
\end{document}
